\documentclass{leveldoc}
\DocUses{foundations}
\DocFingerprint{b3addd0fb5cfa9c8}

\begin{document}

\section{What the level asks}\label{sec:asks}

\leveldescription

One positive ion is launched along $+x$ from the left edge. A metal sphere sits below
the beam and is held at a positive potential, so from outside it repels the ion as a
charge at its centre would. The sphere is a conductor, though: a charge placed near
it pulls the sphere's charge over to the near side. The player may add up to three
charges, of either sign, from the listed magnitudes. There are no magnets.

The ion must enter the detector. There is no acceptance window, no kinetic window
and no gate. A flight ends when the ion touches a charge or the sphere, leaves the
world, or enters the detector (\xsecref{foundations}{sec:board}). The sphere's
obstacle is the metal plus a contact shell of thickness $\val{contact}$.

The reference adds one positive charge under the launch. The sphere lifts the beam,
and that charge lifts it the rest of the way onto the detector's lower edge. This
document builds the field by images, which is exact for one sphere, and integrates
the flight by the classical Runge--Kutta method.

\section{Setup}\label{sec:setup}

\setuptable

The speed of light is $c = \val{cLight}$. The ion has charge $q = \val{qPart}$, mass
$m = \val{mPart}$ and radius $0$, and it is launched with kinetic energy
$T_0 = \val{T0}$. A fixed charge has radius $R = \val{Rcharge}$. The sphere has
centre $(\val{sphereX}, \val{sphereY})$, radius $a = \val{sphereA}$ and potential
$V = \val{sphereV}$. The reference charge is $Q = \val{Qref}$ at
$(\val{xQ}, \val{yQ})$.

The launch is at $(\val{xLaunch}, \val{yLaunch})$, straight along $+x$. Then
$\gamma_0 = \val{gamma0}$, $\abs{\vect{p}_0} = \val{p0}$ and
$v_0 = \val{v0} = \val{vOverC}\,c$ (\xeqref{foundations}{eq:launch}). The detector
is $\val{xDetLo} \le x \le \val{xDetHi}$, $\val{yDetLo} \le y \le \val{yDetHi}$.
The world is $-1 \le x \le \val{xWorld}$, $-1 \le y \le \val{yWorld}$.

\begin{figure}[ht]
\centering
\begin{tikzpicture}
\begin{axis}[axis equal image, width=0.92\linewidth,
  xmin=-1.5, xmax=32.5, ymin=-2, ymax=22.5,
  xlabel={$x$}, ylabel={$y$}, tick label style={font=\small}]
  \draw[black!40, dashed] (-1,-1) rectangle (\rawval{xWorld},\rawval{yWorld});
  \draw[black!25] (0,0) rectangle (30,20);
  \fill[green!50!black, fill opacity=0.18]
    (\rawval{xDetLo},\rawval{yDetLo}) rectangle (\rawval{xDetHi},\rawval{yDetHi});
  \draw[green!40!black]
    (\rawval{xDetLo},\rawval{yDetLo}) rectangle (\rawval{xDetHi},\rawval{yDetHi});
  \fill[blue!30, fill opacity=0.35] (\rawval{sphereX},\rawval{sphereY}) circle[radius=\rawval{sphereA}];
  \draw[blue!50!black] (\rawval{sphereX},\rawval{sphereY}) circle[radius=\rawval{sphereA}];
  \fill[red!70!black] (\rawval{xQ},\rawval{yQ}) circle[radius=\rawval{Rcharge}];
  \addplot[orange!80!black] table {pathdome.dat};
  \addplot[very thick, blue!60!black] table {pathref.dat};
  \draw[-{Stealth[length=5pt]}, thick] (\rawval{xLaunch},\rawval{yLaunch}) -- (\rawval{arrowX},\rawval{arrowY});
  \fill (\rawval{xLaunch},\rawval{yLaunch}) circle[radius=0.08];
\end{axis}
\end{tikzpicture}
\caption{The dome and the two flights. The circle is the sphere and the disc is the
reference charge; the green box is the detector. Blue is the reference flight: it
meets the bottom of the detector. Orange is the same launch with that charge
removed: it leaves through the right edge of the world, below the detector. Both
paths stay above the sphere.}
\label{fig:board}
\end{figure}

\section{Principles and model}\label{sec:model}

\begin{principles}
  \item The field of a point charge and the force on a non-polarizable particle,
    \xeqref{foundations}{eq:coulomb} and \xeqref{foundations}{eq:sphereforce}.
  \item The relativistic equation of motion \xeqref{foundations}{eq:eom}, the launch
    \xeqref{foundations}{eq:launch}, and energy in a static field
    \xeqref{foundations}{eq:energy}. The motion stays in the plane
    (\xsecref{foundations}{sec:plane}).
  \item A metal sphere is an equipotential. Its relaxation is instantaneous on the
    flight's time scale (\physref{2.6}). The field outside one sphere is the image
    field derived in \S\ref{sec:derivation}.
  \item The board's rules (\xsecref{foundations}{sec:board}): obstacles first, then
    the world, then the detector. The classical Runge--Kutta method and the cubic
    Hermite location of an event (\xsecref{foundations}{sec:numerics}).
  \item The radiated power \xeqref{foundations}{eq:larmor}, summed by Simpson's rule
    \xeqref{foundations}{eq:simpson}. It is not put back into the force.
\end{principles}

Outside the sphere the electric field is the Coulomb field of the real charges, of
their images and of one charge $Q_V$ at the centre, plus the image of the ion
itself. There is no magnetic field. The state $(\vect{x}, \vect{p})$ obeys
\xeqref{foundations}{eq:eom} with that $\vect{E}$, from the launch until the first
ending event.

\begin{assumption}{Electrostatic images}\label{as:static}
The metal answers at once. That is exact at infinite $c$, and for finite $c$ it
drops the retardation across the board. The launch speed is $\val{vOverC}\,c$.
The game uses the same static images (\physref{2.6}), so this is the model being
compared, not a correction on top of it.
\end{assumption}

\begin{assumption}{No radiation reaction}\label{as:radiation}
The ion radiates \xeqref{foundations}{eq:larmor}, and that energy is not taken
from the flight. The larger fraction of the launch energy on the two flights here
is $\val{radFracMax}$.
\end{assumption}

\begin{assumption}{The ion's own image}\label{as:self}
The ion induces its own image \eqref{eq:self}. At the closest approach on these
two flights, a centre distance $\val{rNear}$, that force has size $\val{selfForce}$:
a fraction $\val{selfBound}$ of $T_0$ per cell. It is kept in the integration.
The flight never comes near the shell. The smallest clearance of the surface is
$\val{clearMin}$, against a shell of thickness $\val{contact}$ and a sample span
$\val{sampleArc}$.
\end{assumption}

\begin{assumption}{The fixed step}\label{as:step}
The reported step is $h = \val{H}$. Halving it estimates the largest error of an
end state as $\val{errMax}$ (\xsecref{foundations}{sec:numerics}). Ten times that
estimate lies below $10^{-9}$, so the tolerance on a state is the floor
$10^{-9}$.
\end{assumption}

\section{Derivation}\label{sec:derivation}

\subsection{The grounded image}\label{sec:kelvin}

Take the centre as the origin for this paragraph, and put a charge $q$ at distance
$d > a$ along a unit vector $\uvec{u}$. The image is
\begin{equation}\label{eq:kelvin}
  q' = -q\,\frac{a}{d}, \qquad \vect{x}' = \uvec{u}\,\frac{a^2}{d}.
\end{equation}
On the sphere, $\abs{\vect{x}} = a$. Write a surface point as $a\uvec{n}$. Then
\begin{align*}
  \abs{a\uvec{n} - d\uvec{u}}^2 &= a^2 + d^2 - 2ad\,\uvec{n}\cdot\uvec{u}, \\
  \abs{a\uvec{n} - \uvec{u}\,a^2/d}^2
    &= \frac{a^2}{d^2}\bigl(d^2 + a^2 - 2ad\,\uvec{n}\cdot\uvec{u}\bigr).
\end{align*}
So the distance $R$ from the surface point to $q$ and the distance $R'$ to the
image satisfy $R' = (a/d) R$. The potential of the pair on the surface is
\[
  \frac{q}{R} + \frac{q'}{R'}
  = \frac{q}{R} + \frac{-qa/d}{(a/d)R} = 0.
\]
Translated back, a charge at $\vect{p}$ is imaged at
$\vect{c} + (\vect{p}-\vect{c})\,a^2/\abs{\vect{p}-\vect{c}}^2$. With one sphere
there is nowhere further to image that charge: the construction is exact, not a
truncated series (\physref{2.6}).

\subsection{Holding the sphere at a potential}\label{sec:centre}

The pair $(q, q')$ leaves the sphere grounded. A charge $Q_V$ at the centre gives
the constant potential $Q_V/a$ on the surface and does not disturb that
cancellation, so the potential is $V$ when
\begin{equation}\label{eq:centre}
  Q_V = Va.
\end{equation}
Here $Q_V = \val{centreQ}$. The real charges and this centre charge, together
with the images of the real charges, are fixed during the flight. Call their
potential $\phi_\mathrm{ext}$.

The reference charge lies at distance $\val{sourceR}$ from the centre, so its
image is $\val{imageQ}$ at $(\val{imageX}, \val{imageY})$. On sixteen points of
the surface the potential of this fixed set differs from $V$ by at most
$\val{surfRes}$, a fraction $\val{surfRel}$ of $V$: roundoff of the cancellation
above.

\subsection{The ion's own image}\label{sec:self}

The ion at distance $r$ from the centre induces the same grounded image, and the
supply that holds the potential does not add a second centre charge that follows
the ion. The image charge and its place are \eqref{eq:kelvin} with this $q$ and
$d = r$. At the ion, the distance to the image is $(r^2 - a^2)/r$, so the
image's potential there is
\begin{equation}\label{eq:self}
  \phi_\mathrm{self} = -\frac{qa}{r^2 - a^2}.
\end{equation}
The force on the ion is $q\vect{E}$, with $\vect{E}$ the Coulomb field of every
charge except the ion itself, including this image. The image charge is
proportional to $q$ and moves when the ion moves. The work of assembling that
pair is half the final product, and the radial derivative of
$\tfrac{1}{2} q\phi_\mathrm{self}$ reproduces the Coulomb force of the image
(the same check as \xeqref{foundations}{eq:energy} for a fixed source). The
conserved energy is therefore
\begin{equation}\label{eq:image-energy}
  W = (\gamma - 1)mc^2 + q\,\phi_\mathrm{ext} + \tfrac{1}{2} q\,\phi_\mathrm{self}.
\end{equation}
At launch the external part is $q\phi_\mathrm{ext} = \val{qPhi0}$ and the image
part is $\val{halfSelf0}$. The image of the ion at launch is only
$\val{selfQ}$, because the launch is far from the metal.

\subsection{The flight}\label{sec:flight}

There is no closed orbit: the centre charge, the reference charge and its image
are three centres, and the ion's image moves. The state is advanced by
\xeqref{foundations}{eq:rk4} with the fixed step $h = \val{H}$. After each step
the cubic Hermite interpolant of the two ends is sampled at sixty-four fractions
of the step. The first sample that has touched an obstacle, left the world or
entered the detector is refined by bisection, eighty halvings
(\xsecref{foundations}{sec:numerics}). Obstacles are tested before the world, and
the world before the detector. A face coordinate is assigned when the residual
is below $10^{-6}$.

The error of the fine step is $\abs{y_h - y_{2h}}/15$. The same rule estimates
the Simpson error of the radiated energy. Both flights give the same outcome at
$h$ and at $2h$.

\section{The reference solution by hand}\label{sec:hand}

\begin{procedure}
  \item\label{st:launch} \emph{The launch} \xeqref{foundations}{eq:launch}.
    $\gamma_0 = \val{gamma0}$, $\abs{\vect{p}_0} = \val{p0}$, so
    $\vect{p}_0 = (\val{p0}, 0)$ and $v_0 = \val{v0}$.
  \item\label{st:centre} \emph{The centre charge} \eqref{eq:centre}:
    $Q_V = Va = \val{centreQ}$ at $(\val{sphereX}, \val{sphereY})$.
  \item\label{st:image} \emph{The image of the reference charge} \eqref{eq:kelvin}.
    The charge sits at distance $\val{sourceR}$, so
    $q' = -Qa/d = \val{imageQ}$ at $(\val{imageX}, \val{imageY})$.
    The fixed field is this image, the centre charge and the reference charge.
  \item\label{st:march} \emph{The march.} Advance \eqref{eq:image-energy} and
    \xeqref{foundations}{eq:eom} by \xeqref{foundations}{eq:rk4} with step
    $\val{H}$, and locate the ending face as in \S\ref{sec:flight}. The state
    every five time units:
    \begin{center}\small\val{states}\end{center}
  \item\label{st:end} \emph{The arrival.} The ion meets the \val{RefFace} face of
    the detector at $(\val{RefX}, \val{RefY})$, at $t = \val{RefT}$, with
    momentum $(\val{RefPx}, \val{RefPy})$, $\val{RefAng}^\circ$ above $+x$.
    The ordinate is the detector's lower edge. The run takes $\val{RefSteps}$
    steps.
  \item\label{st:check} \emph{The checks.} The fixed potential on the sphere is
    off by at most $\val{surfRes}$. Along the flight, $W$ drifts by at most
    $\val{RefWdev}$. The step-halving estimate of the end state is
    $\val{RefErr}$. The closest centre distance is $\val{RefRclose}$, so the
    clearance of the surface is $\val{RefClear}$.
  \item\label{st:rad} \emph{The radiation.} Simpson's rule on steps of $h/2$
    gives $\val{RefRad}$, a fraction $\val{RefRadFrac}$ of $T_0$. Halving that
    step again changes the integral by the relative amount $\val{RefRadErr}$.
\end{procedure}

\begin{result}[The reference flight]
The ion enters the detector through its lower edge at
$(\val{RefX}, \val{RefY})$ at $t = \val{RefT}$, with momentum
$(\val{RefPx}, \val{RefPy})$. The level is solved.
\end{result}

The march is a few thousand steps of four evaluations. Each evaluation is the
Coulomb sum of four charges: the reference charge, its image, the centre charge
and the ion's image. Nothing else is required.

\section{Comparison with the game}\label{sec:comparison}

\comparisontable

Nine comparisons are in the table: the reference outcome, time and end position,
the radiated energy of the reference, and the outcome, time and end position of
the dome with no added charge. The tolerance on a state is $\val{tolFlight}$.
Ten times the largest step-halving estimate $\val{errMax}$ lies below $10^{-9}$,
so the tolerance is that floor. The radiation tolerance is $10^{-4}$, relative.
The state tolerance and the radiation tolerance were fixed from these estimates
before the comparison.

The largest absolute difference is the reference abscissa,
$\val{cmpRefXDiff}$. The next is the reference time, $\val{cmpRefTDiff}$. Our
truncation $\val{errMax}$ is smaller than either. Both differences sit under the
verified-flight accuracy of about $10^{-10}$
(\xsecref{foundations}{sec:verification}) and well under the tolerance.

The reference ordinate and the dome-only abscissa agree exactly: each is a face
the flight meets, the detector's lower edge and the world's right edge. The
quadrature uncertainty of the radiated energy is $\val{RefRadErr}$. The relative
difference from the game is $\val{cmpRefRadDiff}$, far under the radiation
tolerance. The two quadratures are not the same rule; the formula is.

\section{Other solutions}\label{sec:others}

The sphere alone, which is the empty placement in \code{placements.json}, does
not solve the level.

The ion still leaves along $+x$ and the dome still pushes it upward, but the
push is not enough. The flight leaves through the \val{DomeFace} edge of the
world at $(\val{DomeX}, \val{DomeY})$, at $t = \val{DomeT}$, heading
$\val{DomeAng}^\circ$. The ordinate stays below the detector, whose lower edge
is $\val{yDetLo}$. The clearance of the sphere is $\val{DomeClear}$, and the
energy drifts by $\val{DomeWdev}$.

The reference charge supplies the missing lift. It is positive and sits under
the launch, so it repels the ion upward while the ion is still on the left side
of the board, and the dome carries that extra height to the detector.

\section{What the game computes}\label{sec:engine}

The engine builds the same images (\physref{2.6}). For several spheres it images
again, to a finite depth, and fits a smooth remainder by fundamental solutions
inside each sphere. This level has one sphere, so the second imaging never
starts and the remainder is the roundoff $\val{surfRel}$ measured in
\S\ref{sec:centre}. The ion's own image is included, and the energy diagnostic
is \eqref{eq:image-energy}. The obstacle is the sphere plus the contact shell
$\val{contact}$.

The flight itself is the adaptive eighth-order Runge--Kutta method of
\physref{5} and \physref{6}, verified as in \physref{7}. A verified end state is
good to about $10^{-10}$. The differences in \S\ref{sec:comparison} are inside
that. Radiation is integrated and not applied to the momentum
(\xsecref{foundations}{sec:radiation}, \physref{8}).

\section{Further reading}\label{sec:reading}

The image of a point charge in a grounded sphere, and the centre charge that
raises a sphere to a given potential, are derived in \S\ref{sec:kelvin} and
\S\ref{sec:centre}. The flight uses the Foundations' Runge--Kutta method,
Hermite event location and Simpson rule (\xsecref{foundations}{sec:numerics}).
The engine's conductor model is \physref{2.6}.

\levelfileappendix
\end{document}
